\begin{frame}{LeNet-5에서 중요한 관찰: ``펼쳐서 FC에 넣기''}
  \begin{itemize}
    \item 구조의 ``형''보다 \textbf{완전연결층이 입력을 어떻게
      받아들이는가}가 중요 --- F6은 C5의 특징지도
      (1$\times$1$\times$120, 즉 120개 숫자)를 \textbf{펼쳐서}
      입력으로 씀
    \item LeNet-5는 C5에서 공간 크기를 이미 1$\times$1로 줄여놨기
      때문에 F6 입력이 120으로 \textbf{작지만} ---
    \item \kb{풀링을 덜 하고 채널을 크게 늘린 네트워크}에서는
      이 ``펼쳐진 벡터''가 \textbf{수만$\sim$수십만 길이}가 되어
      \textbf{완전연결층이 파라미터의 대부분}을 차지
  \end{itemize}
  \vskip0.5em
  이 ``펼쳐서 FC에 넣기''가 비효율이라는 인식이
  10.3절 \textbf{전이학습}(``앞쪽 합성곱은 그대로 쓰고 마지막
  FC 헤드만 재학습'')과 \textbf{전역 평균 풀링}의 동기가 된다.
  \vskip0.3em
  공간 32$\times$32 $\to$ 1$\times$1, 채널 1 $\to$ 120 ---
  ``좁아지고 깊어진다'' 경향의 \textbf{원형}.
\end{frame}
